%=====================================================================
% 第 3 节  余二维缺陷与面算符
% 原文：Section 3. Defects of co-dimension two and surface operators
%       （arXiv:1702.06499v2）
%=====================================================================
\section{余二维缺陷与面算符}\label{codim2}

\begin{keytakeaways}
\begin{itemize}
  \item \textbf{物理图像：}把六维 $(2,0)$ 理论看作 M 理论中一叠 M5-膜上的有效理论之后，再加入一叠与原 M5-膜共享部分支撑的 M5'-膜，就在 $\CS$ 类理论中得到余二维面算符（面算符即余二维缺陷）。在二维 sigma 模型中，这类面算符由“零 Nahm 极点”边界条件——即固定 $C$ 上的全纯丛 $\CE$——来描述，对应的膜是 Hitchin 第二纤维化的纤维。
  \item \textbf{数学工具：}Gaiotto--Witten 对 1/2 BPS 边界条件的分类、Nahm 极点渐近行为、非阿贝尔 Hodge（NAH）对应、Hitchin 第二纤维化、仿射李代数共形块及其上的 $\mathcal{A}_{dR}^{\hbar}$/$\mathcal{A}_{B}^{\hbar}$ 模结构、$\bar\partial$-上同调，以及 Virasoro 单值化与 Narasimhan--Seshadri 定理的类比。
  \item \textbf{关键结论：}由固定 $\CE$ 定义的 Lagrangian 膜所给出的空间 $\CH^{\sst (2)}_x$，应与“缓增”（tempered）共形块 $\mathcal{CB}_{\rm te}(C,\hfg_k,\mathcal{E}_x)$ 同构，而非与通常代数定义的 $\mathcal{CB}_{\rm al}(C,\hfg_k,\mathcal{E}_x)$ 同构；后者太大。$\CH^{\sst (2)}_x$ 上存在两个模结构，是支持此猜想的第一条证据。
  \item \textbf{遗留问题：}$\mathcal{CB}_{\rm te}$ 的严格定义在本文中未能给出，作者声明留待另文；Virasoro 情形与 Kac--Moody 情形的讨论都只是“同构于共形块的子空间”的类比论证。
\end{itemize}
\end{keytakeaways}

\bigskip
\noindent\rule{\textwidth}{0.6pt}
\bigskip

M 理论提示了上述框架的一个自然推广。把六维 $(2,0)$ 理论视为 M 理论中一叠 M5-膜上的有效理论，则自然会引入一叠额外的 M5'-膜，使之与原 M5-膜共享一部分支撑（support），如下表所示：
\begin{center}
\begin{tabular}{l || c|c|c|c|c|c|c|c|c|c|c}
膜 & 0 & 1 & 2 & 3 & 4 & 5 & 6 & 7 & 8 & 9 & 10 \\ \hline\hline
M5  & x & x & x & x &   &   & x &   &   &   & x  \\
M5' & x & x &   &   &   &   & x &   & x & x & x \\
\end{tabular}
\end{center}
视所考虑的紧化类型不同，所得对象或是 $\CS$ 类理论中的面算符，或是 $N=4$ SYM 理论中的线算符。这将是第一种修正，而它后来被证明对本文的故事至关重要。

\subsection{$\CS$ 类理论中的面算符}

目前已积累了相当多的证据，表明把 M5-M5'-膜系统在 $C$ 上紧化所得的有效描述，正是被某一类面算符所修正的 $\CS$ 类理论。相关的面算符可通过规定规范场沿二维子流形 $\CM^2\subset\CM^4$ 的某种奇异性来刻画；当 $C$ 的亏格为 $g$、有 $n$ 个穿刺点时，这类奇异性由一组参数 $x=(x_1,\dots,x_{3g-3+n})$ 标定 \cite{AT,FGT}。在存在这类面算符时，可以定义瞬子配分函数的一个推广 $\CZ^{\rm\sst inst}_\si(a;x,\tau;\ep_1,\ep_2)$：除了依赖于标量零模的取值 $a$、与 $C$ 的复结构模 $\tau$ 相联系的紫外规范耦合、以及与 $\Om$-形变相联系的参数 $\ep_1,\ep_2$ 之外，它还额外地依赖于参数 $x$。下标 $\si$ 指决定所用 Lagrangian 表示的裤分解（pants decomposition）。

\cite{Br,AT,KPPW,Ne,Na} 中的显式计算为 AGT 对应的一个推广提供了证据：它把 $\CZ^{\rm\sst inst}_{\si}(a;x,\tau;\ep_1,\ep_2)$ 等同于与仿射李代数 $\hat{\mathfrak{g}}_k$ 的共形块相联系的“手征配分函数” $\CZ^{\rm\sst WZW}_{\si}(a;x,\tau;k)$，该代数的能级为 $k=-h^{\vee}-\ep_2/\ep_1$，而这些共形块由拼接构造（gluing construction）定义。在这一对应下，参数 $x$ 被等同于 $C$ 上全纯 $G$-丛的模空间 $\mathrm{Bun}_G$ 上的坐标。AGT 对应这一推广的证明，将可由 \cite{N15} 所宣布的结果推出，即 $\CZ^{\rm\sst inst}_\si$ 满足 KZB 方程。\footnote{
另见 {\cite{Ne}}，其中有这一方向上的一个较早结果。}

本节的目标是把 \cite{NW,GW} 的进路推广到带余二维面算符的 $\CS$ 类理论。我们将提出以 $\CM_H(G)$ 为目标流形的二维 sigma 模型作为其有效描述，其中 opers 膜 $\mathsf{B}_{\rm op}$ 被一族 Lagrangian 膜 $\mathsf{L}_{x}^{\sst (2)}$ 取代，后者支撑在 Hitchin 第二纤维化（Hitchin's second fibration）位于丛 $\mathcal{E}_x$ 之上的纤维上，而 $\mathcal{E}_x$ 由 $\mathrm{Bun}_G$ 的坐标 $x$ 标记。这些膜 $\mathsf{L}_{x}^{\sst (2)}$ 在此语境中的意义、以及它们与共形块的关系，最早由 E. Frenkel 在 \cite{Fr10} 中提出；这一点曾在 2012 年被指出给作者之一（J.T.）。下文的目标是为这一提议提供进一步的证据。

\subsection{Nahm 极点边界条件}\label{S:Nahm}

Gaiotto 与 Witten 在 \cite{GW1} 中用数据 $(\rho,H,T)$ 对 $N=4$ SYM 的 1/2 BPS 边界条件作了分类，其中 $\rho:\mathfrak{sl}_2\ra \mathfrak{g}$ 是 $\mathfrak{sl}_2$ 到规范群 $G_c$ 的李代数 $\mathfrak{g}$ 中的一个嵌入，$H$ 是 $\rho$ 的像在 $G_c$ 中的中心化子（commutant）的一个子群，而 $T$ 是一个具有 $N=4$ 超对称、且至少具有 $H$ 整体对称性的三维 SCFT。我们只需要其中最简单的两个边界条件。下面先简述 \cite{GW} 中所研究的所谓 Nahm 极点边界条件，它对应于三元组 $(\rho,\mathrm{Id},T_\emptyset)$，其中 $\rho$ 是一个\emph{主}（principal）$\mathfrak{sl}_2$-嵌入，$T_\emptyset$ 表示平凡的三维 SCFT。随后我们将讨论一个更简单的情形：把 $\rho$ 换成将 $\mathfrak{sl}_2$ 映到 $0\in\mathfrak{g}$ 的平凡嵌入；这一情形对我们来说特别重要。

对我们的目的而言，只需在刻画 $N=4$ SYM 中保持某些超对称的场构型的 BPS 方程 \cite{KW} 的解上描述 Nahm 极点边界条件。若把注意力限制在 $\mathbb{R}\times \mathbb{R}_+ \times C$ 上、沿 $\mathbb{R}$ 的平移不变的 BPS 方程解，则得到如下形式的微分方程组
\begin{subequations}\label{BPS}
\begin{align}
&[\,\CD_z\,,\,\CD_{\bar{z}}\,]=0\,,\quad [\,\CD_y\,,\,\CD_{{z}}\,]=0\,,
\quad [\,\CD_y\,,\,\CD_{\bar{z}}\,]=0\,,\label{Fterm}\\
&\sum_{i=1}^3[\,\CD_i\,,\,\CD_i^\dagger\,]=0\,. \label{Dterm}\,,
\end{align}
\end{subequations}
其中已使用记号 $z=x_2+ix_3$ 与 $y=x_1$，而微分算子 $\CD_i$ 取形式\footnote{我们的约定与 \cite{GW} 略有不同。}
\begin{equation}
\begin{aligned}
&\CD_z=\zeta\pa_z+\mathcal{A}_z\,,\\
&\CD_{\bar{z}}=\pa_{\bar{z}}+\mathcal{A}_{\bar{z}}\,,
\end{aligned}
\qquad
\begin{aligned}
&\mathcal{A}_z=\zeta A_z+\phi_z\,,\\
&\mathcal{A}_{\bar{z}}=A_{\bar{z}}+\zeta\phi_{\bar{z}}\,,
\end{aligned}
\qquad \CD_y=\pa_{y}+A_y-i\phi_y\,.
\end{equation}
参数 $\zeta$ 决定了哪些超对称被保持。\cite{GW} 提出：\rf{BPS} 的解空间模紧致规范变换，同构于“F-项”方程 \rf{Fterm} 的解空间模复规范变换所得的模空间。方程 $[\CD_z,\CD_{\bar{z}}]=0$ 在每个固定的 $y$ 处确定 $C$ 上的一个平坦复联络。\rf{Fterm} 中其余的方程则表明，这一平坦联络对 $y$ 的依赖由一个复规范变换来实现。

Nahm 极点型的边界条件在 \cite{GW} 中是这样定义的：要求 \rf{BPS} 的解具有如下形式的奇异性
\begin{equation}\label{Nahm}
\mathcal{A}_z\,\underset{y\ra 0}{\sim}\,\mathfrak{t}_- \,y^{-1}+\CO(y^0)\,,\qquad
\mathcal{A}_{\bar{z}}\,\underset{y\ra 0}{\sim}\,\CO(y^0)\,,\qquad
\mathcal{A}_1\,\underset{y\ra 0}{\sim}\,\mathfrak{t}_3 \,y^{-1}+\CO(y^0)\,,
\end{equation}
其中 $\mathfrak{t}_+=\mathfrak{t}_1+i\mathfrak{t}_2$，而 $\mathfrak{t}_i$（$i=1,2,3$）是 $\mathfrak{g}$ 的某个主 $\mathfrak{sl}_2$ 子代数的生成元。总可以通过一个规范变换把 $\mathcal{A}_{\bar{z}}$ 取为零，于是我们在每个 $y$ 处得到的 $C$ 上的平坦联络可表示为一个局部系统 $(\mathcal{E}_{y},\nabla_{y}')$，它由一个全纯丛和一个全纯联络 $\nabla_{y}'=dz(\pa_z+\mathcal{A}_z(z;y))$ 组成。在 $\mathfrak{g}=\mathfrak{sl}_2$ 的情形，我们可以把 \rf{Nahm} 的第一个条件改述为如下条件：存在一组截面基 $s=\{s_1,s_2\}$，使得 $\mathcal{A}$ 相对于它具有形式 $\mathcal{A}=g\tilde{\mathcal{A}}g^{-1}+gd g^{-1}$，其中
\begin{equation}\label{Nahm'}
\tilde{\mathcal{A}}_z\,\underset{y\ra 0}{\sim}\,\bigg(\begin{matrix} 0 & t\\ 1 & 0\end{matrix}\bigg)+\CO(y^1) \,,
\qquad g\underset{y\ra 0}{\sim} \bigg(\begin{matrix} y^{1/2}& 0 \\  0& y^{-1/2} \end{matrix}\bigg) +\CO(y^0).
\end{equation}
此时水平截面 $(d+\mathcal{A})s=0$ 的第一个分量 $s_1$ 将按 $y^{1/2}$ 的方式消失。正如 \cite{GW} 所解释的，这意味着 $C$ 上的局部系统 $\lim_{y\ra 0}^{}(\tilde{\mathcal{E}}_{y},\tilde{\nabla}_{y}')$ 是一个 oper。

Nahm 极点边界条件的特点是：由于主 $\mathfrak{sl}_2$-嵌入的中心化子是平凡的，它把 $G_c$ 极大地破缺。与之相反的另一个极端，则对应于平凡的 $\mathfrak{sl}_2$-嵌入：通过固定规范场 $\mathcal{A}_{\bar{z}}$ 的边界值，得到与三元组 $(0,\mathrm{Id},T_\emptyset)$ 相联系的类似边界条件。

在约化到二维时，在 $I$ 的边界上固定 $\mathcal{A}_{\bar{z}}$ 意味着在 $C$ 上固定了一个全纯丛，而复规范场的 $(1,0)$ 部分不受约束。偶对 $(\CE,\nabla_\zeta')$ 的模空间 $\CM_{dR}(G)$ 经由非阿贝尔 Hodge（NAH）对应$^{(g)}$ 同构于 Hitchin 模空间 $\CM_H(G)$。因此固定 $\CE$ 就在 $\CM_H(G)$ 中定义了一个子流形，它关于全纯辛形式 $\Omega_\zeta$ 是 Lagrangian 的，并且关于复结构 $I_\zeta$ 是全纯的。当 $\zeta=i$ 时有 $\Omega_\zeta=\Omega_J$、$I_\zeta=J$，于是由零 Nahm 极点边界条件约化所得的膜被认定为 A-模型中（辛结构取 $\omega_I$，即 \cite{NW} 中所用者）的一个 (A,B,A)-膜。

我们提出：零 Nahm 极点边界条件代表着余二维面算符的存在。的确，如 \cite{FGT} 所论证的，余二维面算符的存在会自然地引入额外的背景数据，而这些数据在几何上可以表示为 $C$ 上一个全纯丛的选取。

\subsection{与共形块的关系}

$\CM_H(G)$ 中由固定 $\CE$ 的偶对 $(\CE,\nabla')$ 所定义的 Lagrangian 子流形 $\mathsf{L}_{\CE,\zeta}^{\sst (2)}$，有时被称为 Hitchin 第二纤维化的纤维。取 $\mathrm{Bun}_G$ 的一组坐标 $x$，我们将用 $\mathsf{L}_{x}^{\sst (2)}$ 表示 $\mathsf{L}_{\CE_x,i}^{\sst (2)}$，其中 $\CE_x$ 是代表由坐标 $x$ 所指定的 $\mathrm{Bun}_G$ 中那一点的一个丛。

上面提到的、关于余二维面算符存在时 AGT 对应之推广的结果提示：空间 $\CH^{\sst (2)}_x={\rm Hom}_{\CM_H(G)}(\mathcal{B}_{\rm cc},\mathsf{L}_{x}^{\sst (2)})$ 可以等同于 $C$ 上能级为 $k=-h^{\vee}-\frac{\ep_2}{\ep_1}$ 的仿射李代数 $\hat{\mathfrak{g}}_k$ 的共形块空间。

然而必须指出：在眼下的无穷维情形中，要为 $\CH^{\sst (2)}_x$ 以及相应的共形块空间找到恰当的定义并非平凡。共形块通常的代数定义给出一个空间 $\mathcal{CB}_{\rm al}(C,\hfg_k,\mathcal{E}_x)$，它太大，对我们并不适用。我们需要考虑其中一个子空间，记作 $\mathcal{CB}_{\rm te}(C,\hfg_k,\mathcal{E}_x)$，即“缓增”（tempered）共形块。这两个空间之间的关系在某些方面类似于形式幂级数 $\sum_{n\in\BZ}a_nz^n$ 的空间 $\mathcal{K}_{\rm al}$ 与单位圆上缓增分布的空间 $\mathcal{K}_{\rm te}$ 之间的关系。后者可以在分布意义下用 Fourier 级数 $\sum_{n\in\BZ}a_ne^{in\si}$ 表示，从而给出 $\mathcal{K}_{\rm te}$ 到 $\mathcal{K}_{\rm al}$ 的一个嵌入。然而，“缓增”这一要求对系数 $a_n$ 施加了增长条件，使得 $\mathcal{K}_{\rm te}$ 严格小于 $\mathcal{K}_{\rm al}$。

E. Frenkel 在 \cite{Fr10} 中为 $\CH^{\sst (2)}_x$ 提出了一个代数对应物，此处记作 $\CH^{\sst (2)}_{{\rm al},x}$；他在那里证明了 $\CH^{\sst (2)}_{{\rm al},x}\simeq \mathcal{CB}_{\rm al}(C,\hfg_k,\mathcal{E}_x)$。我们将论证：在此语境下应当为 $\CH^{\sst (2)}_{x}$ 采用另一个不同的定义，从而得到同构 $\CH^{\sst (2)}_{x}\simeq \mathcal{CB}_{\rm te}(C,\hfg_k,\mathcal{E}_x)$，其中 $\mathcal{CB}_{\rm te}(C,\hfg_k,\mathcal{E}_x)$ 是 $\mathcal{CB}_{\rm al}(C,\hfg_k,\mathcal{E}_x)$ 中“缓增”共形块的子空间。在这篇短文里我们无法完整讨论 $\mathcal{CB}_{\rm te}(C,\hfg_k,\mathcal{E}_x)$ 的定义，拟另文再谈。下文将改为讨论来自 $2d$ sigma 模型的、支持 $\CH^{\sst (2)}_{x}\simeq \mathcal{CB}_{\rm te}(C,\hfg_k,\mathcal{E}_x)$ 的证据。

\subsubsection{模结构}

首先注意到，$\CH^{\sst (2)}_x$ 上带有\emph{两个}自然的模结构，它们来自与代数 $\mathcal{A}_{\rm cc}^\hbar={\rm Hom}_{\CM_G(C)}(\mathsf{B}_{\rm cc},\mathsf{B}_{cc})$ 的元素相联系的顶点算符。对于我们所关心的情形，应借助 NAH 对应把 $\CM_H(G)$ 表示为 $C$ 上平坦复联络的模空间；这样便可以考虑分别来自以下两种表示的两种代数结构：把 $\CM_H(G)$ 表示为特征簇$^{(g)}$ $\CM_{B}(G)$，以及表示为偶对 $(\CE,\nabla')$ 的模空间 $\CM_{dR}(G)$。不同的代数结构决定了在 $\mathcal{A}_{\rm cc}$ 中定义顶点算符时所用基本场变量的不同集合。相应地，可以区分
$\mathcal{A}_{B}^\hbar={\rm Hom}_{\CM_B(C)}(\mathsf{B}_{\rm cc},\mathsf{B}_{cc})$
与
$\mathcal{A}_{dR}^\hbar={\rm Hom}_{\CM_{dR}(C)}(\mathsf{B}_{\rm cc},\mathsf{B}_{cc})$，
其中 $\hbar=\ep_1/\ep_2$。

一方面，\cite{NW} 论证了代数 $\mathcal{A}_{B}^\hbar$ 可以等同于 $\CM_{B}(G)$ 上函数的量子化代数。另一方面，\cite{Ka08} 中的论证可以施用于 \cite{NW} 所研究的情形，从而得出结论：$\mathcal{A}_{dR}^\hbar\simeq \CD_{\hbar}$，即 $\CM_{dR}(G)$ 上函数的量子化代数；而后者又可以等同于 $\mathrm{Bun}_G$ 上某个线丛 $\CL^{-h^\vee-\ep_2/\ep_1}$ 上的微分算子代数。

众所周知，能级为 $k=-h^{\vee}-\frac{\ep_2}{\ep_1}$ 的仿射李代数 $\hat{\mathfrak{g}}_k$ 的共形块空间 $\mathcal{CB}_{\rm al}(C,\hat{\mathfrak{g}}_k,\mathcal{E}_x)$ 具有关于 $\mathcal{A}_{dR}^{\hbar}$ 的模结构。这个 $\mathcal{A}_{dR}^{\hbar}$-模结构是其定义性 Ward 恒等式的直接推论。

另一方面，$\mathcal{A}_{B}^{\hbar}$-模结构可以这样在共形块空间上定义：用 $\hat{\mathfrak{g}}_k$ 的退化表示来定义 \cite{AGGTV,DGOT} 中 Verlinde 线算符在此情形的类似物。我们将在别处证明：仿射李代数共形块上 Verlinde 线算符的代数可以等同于 $\mathcal{A}_{B}^{\hbar}$。现有的定义只对用拼接构造得到的共形块有效。这类共形块具有特别好的性质，其中之一是：典范联络可以在整个 $\mathrm{Bun}_G^{\rm vs}$ 上积分，而 $\mathrm{Bun}_G^{\rm vs}$ 是 $\mathrm{Bun}_G$ 中包含“非常稳定”（very stable）丛（即不承认幂零 Higgs 场的丛）的子集；对 $\mathcal{CB}^{\rm al}(C,\hat{\mathfrak{g}}_k,\mathcal{E}_x)$ 的一般元素而言则并非如此。这提示：$\mathcal{A}_{B}^{\hbar}$-模结构只能定义在“缓增”共形块的适当子空间 $\mathcal{CB}_{\rm te}(C,\hat{\mathfrak{g}}_k,\mathcal{E}_x)$ 上。

$\CH^{\rm (2)}_x$ 上存在两个模结构，可以看作支持所猜想同构
$\CH^{\rm (2)}_x\simeq\mathcal{CB}_{\rm te}(C,\hat{\mathfrak{g}}_k,\mathcal{E}_x)$
的第一条证据。下文将给出更多证据。

\subsubsection{与 $\bar\partial$-上同调的关系}

先回到 Nekrasov 与 Witten 在 \cite{NW} 中所研究的情形。首先注意，空间 ${\rm Hom}_{\CM_H(G)}(\mathsf{B}_{\rm cc}, \mathsf{B}_{\rm op})$ 可以在辛结构为 $\omega=\omega_I$ 的 A-模型中得到实现。长条（strip）上的边界条件取为典范余迷向 A-膜 $\mathsf{B}_{\rm cc}$ 与 Lagrangian A-膜 $\mathsf{B}_{\rm op}$。Chan--Paton 曲率 $F$ 由 Hitchin 模空间上的辛形式 $\omega_J$ 给出。于是由典范余迷向 A-膜 $\mathsf{B}_{\rm cc}$ 所确定的复结构为 $\omega^{-1}F=K$。\cite{NW} 论证：在此情形下，空间 ${\rm Hom}_{\CM_H(G)}(\mathsf{B}_{\rm cc}, \mathsf{B}_{\rm op})$ 可以等同于 $\CM_{dR}(G)$ 中 opers 子空间 $\mathrm{Op}_{\fsl_2}(C)$ 上、线丛 $K^{{1}/{2}}_{\rm op}$ 的全纯截面空间。事实上，${\rm Hom}_{\CM_H(G)}(\mathsf{B}_{\rm cc}, \mathsf{B}_{\rm op})$ 将可实现为长条上 A-模型零模的函数空间。对于代表 $\mathrm{Op}_{\fsl_2}(C)$ 上坐标的 sigma 模型场，我们在长条的两端都有 Neumann 型边界条件；而其余的场只在一端有 Dirichlet 型边界条件。因此 A-模型的零模可以用 $\mathrm{Op}_{\fsl_2}(C)$ 上的坐标表示。于是可以援引 \cite{NW} 中简述的标准论证，把长条上 sigma 模型的超对称基态空间等同于取值于 $K^{{1}/{2}}_{\rm op}$ 的 $\bar{\partial}$-上同调。

\subsubsection{与共形块的关系：Virasoro 情形}\label{HomVirConf}

注意 $\mathrm{Op}_{\fsl_2}(C)$（非典范地）同构于复向量空间 $H^0(C,K^2_C)$，后者在拓扑上是平凡的。而空间 $H^0(C,K^2_C)$ 典范同构于 Teichmüller 空间 $\mathcal{T}(C)$ 的余切纤维 $T^*\mathcal{T}(C)|_{C}$。对我们来说重要的是要看到：存在一个从余切空间 $T^*\mathcal{T}(C)|_{C}$ 到 Teichmüller 空间 $\mathcal{T}(C)$ 的映射，它把这两个空间上的自然复结构联系起来。

为了在 $\mathcal{T}(C)$ 上引入复结构，可以用调和 Beltrami 微分来描述 Teichmüller 变分：它们是满足 $\pa_z(\eta\mu)=0$ 的 $(-1,1)$-形式 $\mu$，其中 $\eta$ 是由 $C$ 上复结构唯一确定的双曲度量。于是调和 Beltrami 微分构成的向量空间上的复结构，就定义了 $\mathcal{T}(C)$ 上的复结构。对每个全纯 $(2,0)$-形式 $\theta$，可以配上相应的调和 Beltrami 微分 $\eta^{-1}\bar{\theta}$。这给出一个从 $H^0(C,K^2_C)$ 到调和 Beltrami 微分空间的复反线性映射，它表达了 $\mathcal{T}(C)$ 上与 $H^0(C,K^2_C)$ 上自然复结构之间的关系。由此可以把 $\mathrm{Op}_{\fsl_2}(C)$ 上的全纯函数空间，等同于 $\mathcal{T}(C)$ 上（反）全纯函数的空间 ${\rm Fun}_{\rm hol}(\mathcal{T}(C))$。

另一方面，已知 ${\rm Fun}_{\rm hol}(\mathcal{T}(C))$ 同构于 Virasoro 共形块的（子）空间。闭曲面\footnote{为简化叙述，我们暂时限于这一情形。} $C$ 上的共形块定义为 Virasoro 代数真空表示 $V_0$ 上的线性泛函 $f$，它满足共形 Ward 恒等式，后者刻画 $f$ 在一个点之外全纯的向量场代数之自然作用下的不变性。共形块构成的向量空间带有一个典范的射影平坦联络，它由 Virasoro 代数在 $V_0$ 上的作用定义。这一典范联络的曲率可以在复结构模空间 $\CM_{g}$ 上局部地平凡化。在所有代数定义的共形块空间的适当子空间上，可以对所得的平坦联络积分，从而在带局部坐标 $\tau$ 的 $\CM_{g}$ 的开子集上定义水平截面 $f_{\tau}$。共形块在最高权向量 $v_0$ 上的取值 $\CZ_f(\tau)=f_\tau(v_0)$ 称为手征配分函数。Virasoro 单值化定理与共形 Ward 恒等式把 $\CZ_f(\tau)$ 的导数与刻画共形块 $f$ 的数据联系起来。由此得到（局部定义的）函数 $\CZ(\tau)$ 与“可积”共形块 $f$ 之间的一一对应。这一对应本质上是典范的：Virasoro 单值化描述了 $\CM_{g}$ 的局部结构，而这一点已经编码在共形块的定义之中。

\subsubsection{与共形块的关系：Kac--Moody 情形}

可以论证 $\CH_x^{\sst (2)}$ 同构于 $\hat{\mathfrak{g}}_k$ 的、被丛 $\CE_x$ 扭（twisted）的共形块空间 $\mathcal{CB}(C,\hat{\mathfrak{g}}_k,\CE_x)$，论证方式与第 \ref{HomVirConf} 节类似。依照那里所述的同样推理，将会得到结论：$\CH_{\mathcal{E}}$ 同构于 $H^0(C,{\rm End}(\mathcal{E})\otimes K_C)$ 上的全纯函数空间。空间 $H^0(C,{\rm End}(\mathcal{E})\otimes K_C)$ 正是全纯余切空间 $T^*{\rm Bun}_G$。$T{\rm Bun}_G$ 上存在一个典范的 Hermite 度量，它类似于 Teichmüller 理论中的 Weil--Petersson 度量，其定义以 Narasimhan--Seshadri 定理作为单值化定理的类比。该度量把全纯余切空间等同于反全纯切空间，从而把 $H^0(C,{\rm End}(\mathcal{E})\otimes K_C)$ 与 $T{\rm Bun}_G$ 上的自然复结构联系起来。

$T{\rm Bun}_G$ 上的全纯函数空间自然地嵌入 $\on{Jet}_\mathcal{E}$，即 $\mathcal{E}$ 在 ${\rm Bun}_G$ 中形式邻域上的函数空间。剩下只需注意 $\on{Jet}_\mathcal{E}$ 同构于共形块空间，$\on{Jet}_\mathcal{E}\simeq\mathcal{CB}_{\rm al}(C,\hat{\mathfrak{g}}_k,\CE_x)$ \cite[Section 18.2]{FBZ}。E. Frenkel 正是利用这一观察在 \cite{Fr10} 中证明了我们论断的代数对应物 $\CH^{\sst (2)}_{{\rm al},x} \simeq \mathcal{CB}_{\rm al}(C,\hfg_k,\CE_x)$。由此可见，$T{\rm Bun}_G$ 上的全纯函数空间可以等同于 $\mathcal{CB}(C,\hat{\mathfrak{g}}_k,\CE_x)$ 中“性质良好”（well-behaved）的共形块的子空间，这进一步支持了上面提出的猜想 $\CH^{\sst (2)}_{x} \simeq \mathcal{CB}_{\rm te}(C,\hfg_k,\CE_x)$。

\bigskip
\noindent\rule{\textwidth}{0.6pt}

% [Tutor Note: 式 \eqref{Nahm} 与紧接其后的文字不自洽：正文定义的是 $\mathfrak{t}_+=\mathfrak{t}_1+i\mathfrak{t}_2$，而式中出现的是 $\mathfrak{t}_-$（极点系数）与 $\mathfrak{t}_3$（$A_y$ 的极点系数）。按通常约定 $\mathfrak{t}_+ = \mathfrak{t}_1+i\mathfrak{t}_2$、$\mathfrak{t}_- = \mathfrak{t}_1-i\mathfrak{t}_2$，则 $A_z$ 的极点系数应是 $\mathfrak{t}_-$、而文中只给出了 $\mathfrak{t}_+$ 的定义，读者若自行补 $\mathfrak{t}_-$ 须取复共轭（在 $\mathfrak{sl}_2$ 的实形式下 $\mathfrak{t}_-=\mathfrak{t}_+^\dagger$）。另注意第三条中的 $\mathcal{A}_1$ 与前文 $\CD_y=\pa_y+A_y-i\phi_y$ 中的 $A_y$ 是同一个分量的两种写法（$\mathcal{A}$ 带花体在别处专指 twisted connection），核对 Nahm 极点条件时不要把它误当成第四分量。]
% [Tutor Note: 式 \eqref{BPS} 到式 \eqref{Fterm} 之间被略去的一步是 Hitchin 方程（此处为三维约化后的 Nahm 方程）的超 Kähler 商论证：\eqref{Dterm} 的矩映射方程可通过一个复规范变换解出，而该复规范变换的选取不改变 \eqref{Fterm} 的复平坦性；也就是说，“BPS 解模紧致规范变换”与“F-项解模复规范变换”同构，靠的是 Kempf--Ness 型定理在无穷维的类比（Donaldson--Uhlenbeck--Yau），并且在 $y\ra0$ 处还需对渐近行为加正则性条件，否则模空间不具有限维性。作者引用 \cite{GW} 而未写出此论证。另注意算子定义中 $\zeta$ 的幂次并不对称：$A_z$ 带一个 $\zeta$、$\phi_{\bar z}$ 带一个 $\zeta$，而 $A_{\bar z}$、$\phi_z$ 不带；作者已加脚注声明与 \cite{GW} 约定不同；重标度 $\zeta$ 会同时改变 $\Om_\zeta$ 与 $I_\zeta$，从而改变膜的 A/B 类型判定，与 \cite{NW} 对照时不可混用两套归一化。]
% [Tutor Note: “固定 $\CE$ 的子流形关于 $\Omega_\zeta$ 是 Lagrangian 的”一句是本文后面所有膜判定的基础，但此处只作断言。实际上要用到两件事：(i) NAH 对应给出的是实解析（超 Kähler）同构，而非代数同构，它把 $\CM_{dR}$ 上的复结构 $I_\zeta$ 与全纯辛形式 $\Omega_\zeta$ 搬到 $\CM_H$ 一侧；(ii) Hitchin 第二纤维化 $\CM_H(G)\to\Bun_G$（丢掉 Higgs 场）的纤维维数为 $\dim\Bun_G$，恰为 $\CM_H$ 复维数的一半，且 $\Omega_\zeta$ 沿纤维为零——这一点来自 Hitchin 可积系统中 Poisson 对合函数的构造，需要单独验证，不能由“固定 $\CE$”直接看出。]
% [Tutor Note: $\mathcal{A}_{dR}^\hbar\simeq\CD_\hbar$ 并“等同于 $\CL^{-h^\vee-\ep_2/\ep_1}$ 上微分算子代数”这一步用的是 Beilinson--Drinfeld 关于临界能级中心（其谱即 opers）的定理之量子化。核对幂次：因 $k=-h^\vee-\ep_2/\ep_1$，指数 $-h^\vee-\ep_2/\ep_1$ 恰好就是 $k$，即线丛取 $\CL^{k}$；若按常见的 $\CL^{\otimes(k+h^\vee)}$ 写法，则指数是 $-\ep_2/\ep_1$，两者相差 $h^\vee$，极易写错。另外 $\hbar=\ep_1/\ep_2$ 与 $k+h^\vee=-\ep_2/\ep_1=-1/\hbar$ 互为倒数，后文比较 $W_{1/2\hbar}$ 时须一并核对。]
% [Tutor Note: 关于 $(-1,1)$-形式与条件 $\pa_z(\eta\mu)=0$ 的指标核算：$\mu$ 按 $\mu=\eta^{-1}\bar\theta$ 的类型计为 $K^{-1}\bar K$ 的截面（即 $(-1,1)$-型），与全纯二次微分 $\theta\in H^0(C,K^2_C)$ 配对后 $\mu\theta$ 是 $(1,1)$-形式方可积分；$\eta\mu=\bar\theta$ 的类型是 $(0,2)$，故 $\pa_z(\eta\mu)=0$ 应理解为“$\eta\mu$ 是 $C$ 上的全纯二次微分”（$\pa_z$ 作用在系数函数上），它同时给出 $\mu$ 的调和性。若把 $\mu$ 误当作 $(1,-1)$-型（即反变指标写反），$\pa_z$ 与 $\bar\pa_{\bar z}$ 的位置就会颠倒，$(-1,1)$ 与 $(1,-1)$ 也常被混用。]
% [Tutor Note: 两处易被带过的措辞：(i) “very stable 丛”的严格定义是“不承认\emph{非零}幂零 Higgs 场”，原文漏了“非零”二字（零 Higgs 场总是存在的），这一定义源于 Drinfeld 对 very stable 丛的引入（Laumon 讨论了其变体），且与“典范联络可在 $\mathrm{Bun}_G^{\rm vs}$ 上积分”直接相关；(ii) 缓增空间与形式幂级数空间之间，原文说嵌入方向是 $\mathcal{K}_{\rm te}\hookrightarrow\mathcal{K}_{\rm al}$，这是对的（多项式增长的 Fourier 系数必为形式幂级数），但作者未给出 $\mathcal{CB}_{\rm te}$ 的任何定义，读者不应把“tempered”误解为 $L^2$ 完备化——它对应的是 Fourier 系数（即共形块在穿刺点附近的模参数展开系数）的多项式增长条件。]
